% ============================================================
% Close-up: ScientistsAgent — Workflow Interface View
% Fragment — % ============================================================
% Close-up: ScientistsAgent — Workflow Interface View
% Fragment — % ============================================================
% Close-up: ScientistsAgent — Workflow Interface View
% Fragment — % ============================================================
% Close-up: ScientistsAgent — Workflow Interface View
% Fragment — \input{figs/closeup_scientists}
% ============================================================
\begin{figure}[htbp]
\centering
\resizebox{.78\textwidth}{!}{%
\begin{tikzpicture}[node distance=0.35cm, font=\sffamily,
    stepbar/.style={rectangle, rounded corners=6pt, thick,
        minimum width=20cm, align=center, font=\sffamily,
        drop shadow={opacity=0.10}, inner sep=6pt, text width=19cm},
    rwpanel/.style={rectangle, rounded corners=5pt, thick,
        minimum width=20cm, align=left, inner sep=6pt,
        font=\sffamily\small, text width=18.5cm},
    infopanel/.style={rectangle, rounded corners=5pt, draw=gray!50, thick,
        minimum width=20cm, align=left, fill=white,
        inner sep=6pt, font=\sffamily\small, text width=18.5cm},
    toolpanel/.style={rectangle, rounded corners=5pt, draw=gray!60, thick,
        minimum width=20cm, align=left, fill=gray!4,
        inner sep=6pt, font=\sffamily\small, text width=18.5cm},
    expertcard/.style={rectangle, rounded corners=4pt, draw=gray!50,
        thick, minimum width=4.2cm, minimum height=2.2cm,
        align=center, font=\sffamily\scriptsize, inner sep=5pt},
    flowline/.style={->, >=Stealth, very thick, draw=gray!40},
    fanline/.style={->, >=Stealth, thick, draw=teal!40},
    consumerlabel/.style={font=\sffamily\small\itshape, text=gray!60!black,
        align=center, text width=20cm}
]
    % === STEP BAR ===
    \node[stepbar, draw=teal!50!black, fill=agentTeal] (step)
        {\iconBulb{teal!60!black}\;\;
         \textbf{\large ScientistsAgent --- Multi-Expert Hypothesis
          Synthesis}\\[3pt]
         {\small Planner: \textit{``Synthesise mechanisms via expert
          panel, score hypotheses, propose experiments.''}}\\[2pt]
         {\scriptsize\sffamily\color{gray!60!black}
          Shown as a fourth step here; the planner dynamically
          decides invocation order and may loop back to earlier
          agents after synthesis reveals knowledge gaps.}};

    % === READS ===
    \node[rwpanel, draw=green!40!black, fill=agentGreen,
          below=0.4cm of step] (reads)
        {\textbf{$\blacktriangledown$ Reads from AgentState}
         {\scriptsize (heaviest reader --- consumes all prior output)}\\[5pt]
         $\bullet$\; \texttt{ALL previous\_results} {\scriptsize (str)} ---
          full accumulated output from Steps 1--3
          (omics + KG + literature)\\[2pt]
         $\bullet$\; \texttt{top\_genes} {\scriptsize (list[str])} ---
          ranked DEGs \emph{from OmicMiningAgent}\\[2pt]
         $\bullet$\; \texttt{paper\_dois} {\scriptsize (list[str])} ---
          literature DOIs \emph{from PubMedResearcher},
          used for citation cross-referencing};

    \draw[flowline] (step.south) -- (reads.north);

    % === PROMPT HIGHLIGHTS ===
    \node[infopanel, below=0.4cm of reads] (prompt)
        {\textbf{System Prompt Highlights}
         {\scriptsize (from \texttt{ScientistsAgent.system\_message})}\\[5pt]
         \textbf{1.}\; ``Analyse the incoming research task and
          determine which expert(s) are best suited''\\[2pt]
         \textbf{2.}\; ``Delegate with specific, tailored instructions
          for each expert''\\[2pt]
         \textbf{3.}\; ``Synthesise their responses into a unified,
          high-quality scientific analysis''\\[2pt]
         \textbf{4.}\; ``Prefer experts with \textbf{(RAG KB loaded)}
          --- they can retrieve evidence from real papers''\\[2pt]
         \textbf{5.}\; ``ALWAYS include a bioinformatics perspective
          when data-analysis methodology is involved''\\[2pt]
         \textbf{6.}\; ``ALWAYS include a biostatistics perspective
          when statistical claims or ageing are involved''\\[2pt]
         \textbf{7.}\; ``For hypothesis-generation tasks, use ALL
          relevant experts for multi-perspective coverage''};

    \draw[flowline] (reads.south) -- (prompt.north);

    % === ROUTING + EXPERT PANEL ===
    \node[toolpanel, below=0.4cm of prompt] (routing)
        {\textbf{PI Router:}\;
         Structured LLM output via \texttt{ExpertRouting}
         Pydantic model $\rightarrow$ selects 2--3 experts with
         tailored sub-queries};

    \draw[flowline] (prompt.south) -- (routing.north);

    % --- 4 Expert cards in a row ---
    \node[expertcard, fill=yellow!8,
          below=0.5cm of routing, xshift=-7cm] (exp1)
        {\textbf{Genomics-}\\
         \textbf{Expert}\\[3pt]
         \colorbox{green!15}{\scriptsize RAG KB}\\[3pt]
         {\tiny gene regulation,}\\
         {\tiny variant interp.,}\\
         {\tiny CRISPRi/a/ko,}\\
         {\tiny pharmacogenomics}};

    \node[expertcard, fill=pink!8,
          right=0.5cm of exp1] (exp2)
        {\textbf{Neuroscience-}\\
         \textbf{Expert}\\[3pt]
         \colorbox{green!15}{\scriptsize RAG KB}\\[3pt]
         {\tiny A$\beta$, tau, TREM2,}\\
         {\tiny microglia DAM,}\\
         {\tiny neuroinflammation,}\\
         {\tiny CSF biomarkers}};

    \node[expertcard, fill=lime!8,
          right=0.5cm of exp2] (exp3)
        {\textbf{Longevity/}\\
         \textbf{BiostatsExpert}\\[3pt]
         \colorbox{green!15}{\scriptsize RAG KB}\\[3pt]
         {\tiny ageing hallmarks,}\\
         {\tiny epigenetic clocks,}\\
         {\tiny mTOR/AMPK/sirtuins,}\\
         {\tiny survival analysis}};

    \node[expertcard, fill=cyan!8,
          right=0.5cm of exp3] (exp4)
        {\textbf{Cancer-}\\
         \textbf{Expert}\\[3pt]
         \colorbox{red!15}{\scriptsize LLM-only}\\[3pt]
         {\tiny tumor microenv.,}\\
         {\tiny oncogenesis,}\\
         {\tiny cancer therapeutics,}\\
         {\tiny resistance mech.}};

    % Fan-out arrows from routing to experts
    \draw[fanline] (routing.south) -- ++(0,-0.25) -| (exp1.north);
    \draw[fanline] (routing.south) -- ++(0,-0.25) -| (exp2.north);
    \draw[fanline] (routing.south) -- ++(0,-0.25) -| (exp3.north);
    \draw[fanline] (routing.south) -- ++(0,-0.25) -| (exp4.north);

    % Shared tool below experts
    \node[toolpanel, minimum width=19cm, text width=18cm,
          below=0.5cm of $(exp2.south)!0.5!(exp3.south)$] (ragtool)
        {\textbf{Shared Tool:}\;
         \texttt{scientist\_rag\_tool(query, expert)}\\[3pt]
         {\scriptsize Queries the specified expert's HyperRAG knowledge
          base built from curated author publications.\;
          Uses hypergraph-based retrieval with OpenAI embeddings and
          GPT-4.1 for answer generation.\;
          \textbf{RAG~KB} experts retrieve evidence from real papers;\;
          \textbf{LLM-only} experts reason without a private KB when
          their author KB has not yet been built.}\\[3pt]
         {\scriptsize \textit{See companion illustration for Scientist
          RAG initialisation and retrieval pipeline.}}};

    % Fan-in to shared tool
    \draw[fanline] (exp1.south) -- ++(0,-0.2) -| (ragtool.north west);
    \draw[fanline] (exp2.south) -- ++(0,-0.2) -|
                   ($(ragtool.north)!0.33!(ragtool.north west)$);
    \draw[fanline] (exp3.south) -- ++(0,-0.2) -|
                   ($(ragtool.north)!0.33!(ragtool.north east)$);
    \draw[fanline] (exp4.south) -- ++(0,-0.2) -| (ragtool.north east);

    % === PI SYNTHESIS ===
    \node[infopanel, draw=teal!50, fill=teal!5,
          below=0.4cm of ragtool, minimum width=20cm,
          text width=18.5cm] (synth)
        {\iconRobot{teal!60!black}\;\;
         \textbf{PI Synthesis}\\[3pt]
         {\scriptsize Integrates expert perspectives into unified
          analysis.\; Ranks mechanistic hypotheses by evidence
          strength.\; Proposes specific validation experiments
          per hypothesis.}};

    \draw[flowline] (ragtool.south) -- (synth.north);

    % === WRITES ===
    \node[rwpanel, draw=blue!30!black, fill=stateBlue,
          below=0.4cm of synth, minimum width=20cm,
          text width=18.5cm] (writes)
        {\textbf{$\blacktriangle$ Writes to AgentState}\\[5pt]
         $\bullet$\; \texttt{agent\_outputs[]} {\scriptsize (str)} ---
          multi-expert synthesis including:\\[2pt]
         \hphantom{$\bullet$\;}
          {\scriptsize --- ranked mechanistic hypotheses with
           Gene$\rightarrow$Pathway$\rightarrow$Phenotype chains}\\[1pt]
         \hphantom{$\bullet$\;}
          {\scriptsize --- per-expert domain perspectives and
           evidence summaries}\\[1pt]
         \hphantom{$\bullet$\;}
          {\scriptsize --- proposed validation experiments with
           quantitative readouts}\\[3pt]
         $\bullet$\; \texttt{expert\_routing\{\}} {\scriptsize (dict)} ---
          which experts were selected, why, and their tailored
          sub-queries};

    \draw[flowline] (synth.south) -- (writes.north);

    % === DOWNSTREAM CONSUMERS ===
    \node[consumerlabel, below=0.3cm of writes, text width=20cm]
          (consumers)
        {$\downarrow$\; Consumed by:\;
         \textbf{GoogleSearcher}
          (hypotheses to validate translationally)\; $\cdot$\;
         \textbf{ReportingNode}
          (hypothesis section, validation plans, expert routing log)};

\end{tikzpicture}%
}% end resizebox
\caption{\textbf{ScientistsAgent} --- workflow interface view.
The heaviest reader of \texttt{AgentState}, this compound agent
consumes all prior omics, KG, and literature output.  A PI-level
router uses a structured \texttt{ExpertRouting} Pydantic model to
select 2--3 of the four domain experts configured in
\texttt{configs/experts.yaml} and dispatch tailored sub-queries.
Genomics, Neuroscience, and Longevity/Biostatistics experts are
backed by curated-author HyperRAG KBs (\emph{RAG~KB}); CancerExpert
is configured but its KB is not yet built and is therefore routed
\emph{LLM-only}.  The PI synthesis node integrates expert
perspectives into ranked hypotheses with proposed validation
experiments.}
\label{fig:closeup-scientists}
\end{figure}

% ============================================================
\begin{figure}[htbp]
\centering
\resizebox{.78\textwidth}{!}{%
\begin{tikzpicture}[node distance=0.35cm, font=\sffamily,
    stepbar/.style={rectangle, rounded corners=6pt, thick,
        minimum width=20cm, align=center, font=\sffamily,
        drop shadow={opacity=0.10}, inner sep=6pt, text width=19cm},
    rwpanel/.style={rectangle, rounded corners=5pt, thick,
        minimum width=20cm, align=left, inner sep=6pt,
        font=\sffamily\small, text width=18.5cm},
    infopanel/.style={rectangle, rounded corners=5pt, draw=gray!50, thick,
        minimum width=20cm, align=left, fill=white,
        inner sep=6pt, font=\sffamily\small, text width=18.5cm},
    toolpanel/.style={rectangle, rounded corners=5pt, draw=gray!60, thick,
        minimum width=20cm, align=left, fill=gray!4,
        inner sep=6pt, font=\sffamily\small, text width=18.5cm},
    expertcard/.style={rectangle, rounded corners=4pt, draw=gray!50,
        thick, minimum width=4.2cm, minimum height=2.2cm,
        align=center, font=\sffamily\scriptsize, inner sep=5pt},
    flowline/.style={->, >=Stealth, very thick, draw=gray!40},
    fanline/.style={->, >=Stealth, thick, draw=teal!40},
    consumerlabel/.style={font=\sffamily\small\itshape, text=gray!60!black,
        align=center, text width=20cm}
]
    % === STEP BAR ===
    \node[stepbar, draw=teal!50!black, fill=agentTeal] (step)
        {\iconBulb{teal!60!black}\;\;
         \textbf{\large ScientistsAgent --- Multi-Expert Hypothesis
          Synthesis}\\[3pt]
         {\small Planner: \textit{``Synthesise mechanisms via expert
          panel, score hypotheses, propose experiments.''}}\\[2pt]
         {\scriptsize\sffamily\color{gray!60!black}
          Shown as a fourth step here; the planner dynamically
          decides invocation order and may loop back to earlier
          agents after synthesis reveals knowledge gaps.}};

    % === READS ===
    \node[rwpanel, draw=green!40!black, fill=agentGreen,
          below=0.4cm of step] (reads)
        {\textbf{$\blacktriangledown$ Reads from AgentState}
         {\scriptsize (heaviest reader --- consumes all prior output)}\\[5pt]
         $\bullet$\; \texttt{ALL previous\_results} {\scriptsize (str)} ---
          full accumulated output from Steps 1--3
          (omics + KG + literature)\\[2pt]
         $\bullet$\; \texttt{top\_genes} {\scriptsize (list[str])} ---
          ranked DEGs \emph{from OmicMiningAgent}\\[2pt]
         $\bullet$\; \texttt{paper\_dois} {\scriptsize (list[str])} ---
          literature DOIs \emph{from PubMedResearcher},
          used for citation cross-referencing};

    \draw[flowline] (step.south) -- (reads.north);

    % === PROMPT HIGHLIGHTS ===
    \node[infopanel, below=0.4cm of reads] (prompt)
        {\textbf{System Prompt Highlights}
         {\scriptsize (from \texttt{ScientistsAgent.system\_message})}\\[5pt]
         \textbf{1.}\; ``Analyse the incoming research task and
          determine which expert(s) are best suited''\\[2pt]
         \textbf{2.}\; ``Delegate with specific, tailored instructions
          for each expert''\\[2pt]
         \textbf{3.}\; ``Synthesise their responses into a unified,
          high-quality scientific analysis''\\[2pt]
         \textbf{4.}\; ``Prefer experts with \textbf{(RAG KB loaded)}
          --- they can retrieve evidence from real papers''\\[2pt]
         \textbf{5.}\; ``ALWAYS include a bioinformatics perspective
          when data-analysis methodology is involved''\\[2pt]
         \textbf{6.}\; ``ALWAYS include a biostatistics perspective
          when statistical claims or ageing are involved''\\[2pt]
         \textbf{7.}\; ``For hypothesis-generation tasks, use ALL
          relevant experts for multi-perspective coverage''};

    \draw[flowline] (reads.south) -- (prompt.north);

    % === ROUTING + EXPERT PANEL ===
    \node[toolpanel, below=0.4cm of prompt] (routing)
        {\textbf{PI Router:}\;
         Structured LLM output via \texttt{ExpertRouting}
         Pydantic model $\rightarrow$ selects 2--3 experts with
         tailored sub-queries};

    \draw[flowline] (prompt.south) -- (routing.north);

    % --- 4 Expert cards in a row ---
    \node[expertcard, fill=yellow!8,
          below=0.5cm of routing, xshift=-7cm] (exp1)
        {\textbf{Genomics-}\\
         \textbf{Expert}\\[3pt]
         \colorbox{green!15}{\scriptsize RAG KB}\\[3pt]
         {\tiny gene regulation,}\\
         {\tiny variant interp.,}\\
         {\tiny CRISPRi/a/ko,}\\
         {\tiny pharmacogenomics}};

    \node[expertcard, fill=pink!8,
          right=0.5cm of exp1] (exp2)
        {\textbf{Neuroscience-}\\
         \textbf{Expert}\\[3pt]
         \colorbox{green!15}{\scriptsize RAG KB}\\[3pt]
         {\tiny A$\beta$, tau, TREM2,}\\
         {\tiny microglia DAM,}\\
         {\tiny neuroinflammation,}\\
         {\tiny CSF biomarkers}};

    \node[expertcard, fill=lime!8,
          right=0.5cm of exp2] (exp3)
        {\textbf{Longevity/}\\
         \textbf{BiostatsExpert}\\[3pt]
         \colorbox{green!15}{\scriptsize RAG KB}\\[3pt]
         {\tiny ageing hallmarks,}\\
         {\tiny epigenetic clocks,}\\
         {\tiny mTOR/AMPK/sirtuins,}\\
         {\tiny survival analysis}};

    \node[expertcard, fill=cyan!8,
          right=0.5cm of exp3] (exp4)
        {\textbf{Cancer-}\\
         \textbf{Expert}\\[3pt]
         \colorbox{red!15}{\scriptsize LLM-only}\\[3pt]
         {\tiny tumor microenv.,}\\
         {\tiny oncogenesis,}\\
         {\tiny cancer therapeutics,}\\
         {\tiny resistance mech.}};

    % Fan-out arrows from routing to experts
    \draw[fanline] (routing.south) -- ++(0,-0.25) -| (exp1.north);
    \draw[fanline] (routing.south) -- ++(0,-0.25) -| (exp2.north);
    \draw[fanline] (routing.south) -- ++(0,-0.25) -| (exp3.north);
    \draw[fanline] (routing.south) -- ++(0,-0.25) -| (exp4.north);

    % Shared tool below experts
    \node[toolpanel, minimum width=19cm, text width=18cm,
          below=0.5cm of $(exp2.south)!0.5!(exp3.south)$] (ragtool)
        {\textbf{Shared Tool:}\;
         \texttt{scientist\_rag\_tool(query, expert)}\\[3pt]
         {\scriptsize Queries the specified expert's HyperRAG knowledge
          base built from curated author publications.\;
          Uses hypergraph-based retrieval with OpenAI embeddings and
          GPT-4.1 for answer generation.\;
          \textbf{RAG~KB} experts retrieve evidence from real papers;\;
          \textbf{LLM-only} experts reason without a private KB when
          their author KB has not yet been built.}\\[3pt]
         {\scriptsize \textit{See companion illustration for Scientist
          RAG initialisation and retrieval pipeline.}}};

    % Fan-in to shared tool
    \draw[fanline] (exp1.south) -- ++(0,-0.2) -| (ragtool.north west);
    \draw[fanline] (exp2.south) -- ++(0,-0.2) -|
                   ($(ragtool.north)!0.33!(ragtool.north west)$);
    \draw[fanline] (exp3.south) -- ++(0,-0.2) -|
                   ($(ragtool.north)!0.33!(ragtool.north east)$);
    \draw[fanline] (exp4.south) -- ++(0,-0.2) -| (ragtool.north east);

    % === PI SYNTHESIS ===
    \node[infopanel, draw=teal!50, fill=teal!5,
          below=0.4cm of ragtool, minimum width=20cm,
          text width=18.5cm] (synth)
        {\iconRobot{teal!60!black}\;\;
         \textbf{PI Synthesis}\\[3pt]
         {\scriptsize Integrates expert perspectives into unified
          analysis.\; Ranks mechanistic hypotheses by evidence
          strength.\; Proposes specific validation experiments
          per hypothesis.}};

    \draw[flowline] (ragtool.south) -- (synth.north);

    % === WRITES ===
    \node[rwpanel, draw=blue!30!black, fill=stateBlue,
          below=0.4cm of synth, minimum width=20cm,
          text width=18.5cm] (writes)
        {\textbf{$\blacktriangle$ Writes to AgentState}\\[5pt]
         $\bullet$\; \texttt{agent\_outputs[]} {\scriptsize (str)} ---
          multi-expert synthesis including:\\[2pt]
         \hphantom{$\bullet$\;}
          {\scriptsize --- ranked mechanistic hypotheses with
           Gene$\rightarrow$Pathway$\rightarrow$Phenotype chains}\\[1pt]
         \hphantom{$\bullet$\;}
          {\scriptsize --- per-expert domain perspectives and
           evidence summaries}\\[1pt]
         \hphantom{$\bullet$\;}
          {\scriptsize --- proposed validation experiments with
           quantitative readouts}\\[3pt]
         $\bullet$\; \texttt{expert\_routing\{\}} {\scriptsize (dict)} ---
          which experts were selected, why, and their tailored
          sub-queries};

    \draw[flowline] (synth.south) -- (writes.north);

    % === DOWNSTREAM CONSUMERS ===
    \node[consumerlabel, below=0.3cm of writes, text width=20cm]
          (consumers)
        {$\downarrow$\; Consumed by:\;
         \textbf{GoogleSearcher}
          (hypotheses to validate translationally)\; $\cdot$\;
         \textbf{ReportingNode}
          (hypothesis section, validation plans, expert routing log)};

\end{tikzpicture}%
}% end resizebox
\caption{\textbf{ScientistsAgent} --- workflow interface view.
The heaviest reader of \texttt{AgentState}, this compound agent
consumes all prior omics, KG, and literature output.  A PI-level
router uses a structured \texttt{ExpertRouting} Pydantic model to
select 2--3 of the four domain experts configured in
\texttt{configs/experts.yaml} and dispatch tailored sub-queries.
Genomics, Neuroscience, and Longevity/Biostatistics experts are
backed by curated-author HyperRAG KBs (\emph{RAG~KB}); CancerExpert
is configured but its KB is not yet built and is therefore routed
\emph{LLM-only}.  The PI synthesis node integrates expert
perspectives into ranked hypotheses with proposed validation
experiments.}
\label{fig:closeup-scientists}
\end{figure}

% ============================================================
\begin{figure}[htbp]
\centering
\resizebox{.78\textwidth}{!}{%
\begin{tikzpicture}[node distance=0.35cm, font=\sffamily,
    stepbar/.style={rectangle, rounded corners=6pt, thick,
        minimum width=20cm, align=center, font=\sffamily,
        drop shadow={opacity=0.10}, inner sep=6pt, text width=19cm},
    rwpanel/.style={rectangle, rounded corners=5pt, thick,
        minimum width=20cm, align=left, inner sep=6pt,
        font=\sffamily\small, text width=18.5cm},
    infopanel/.style={rectangle, rounded corners=5pt, draw=gray!50, thick,
        minimum width=20cm, align=left, fill=white,
        inner sep=6pt, font=\sffamily\small, text width=18.5cm},
    toolpanel/.style={rectangle, rounded corners=5pt, draw=gray!60, thick,
        minimum width=20cm, align=left, fill=gray!4,
        inner sep=6pt, font=\sffamily\small, text width=18.5cm},
    expertcard/.style={rectangle, rounded corners=4pt, draw=gray!50,
        thick, minimum width=4.2cm, minimum height=2.2cm,
        align=center, font=\sffamily\scriptsize, inner sep=5pt},
    flowline/.style={->, >=Stealth, very thick, draw=gray!40},
    fanline/.style={->, >=Stealth, thick, draw=teal!40},
    consumerlabel/.style={font=\sffamily\small\itshape, text=gray!60!black,
        align=center, text width=20cm}
]
    % === STEP BAR ===
    \node[stepbar, draw=teal!50!black, fill=agentTeal] (step)
        {\iconBulb{teal!60!black}\;\;
         \textbf{\large ScientistsAgent --- Multi-Expert Hypothesis
          Synthesis}\\[3pt]
         {\small Planner: \textit{``Synthesise mechanisms via expert
          panel, score hypotheses, propose experiments.''}}\\[2pt]
         {\scriptsize\sffamily\color{gray!60!black}
          Shown as a fourth step here; the planner dynamically
          decides invocation order and may loop back to earlier
          agents after synthesis reveals knowledge gaps.}};

    % === READS ===
    \node[rwpanel, draw=green!40!black, fill=agentGreen,
          below=0.4cm of step] (reads)
        {\textbf{$\blacktriangledown$ Reads from AgentState}
         {\scriptsize (heaviest reader --- consumes all prior output)}\\[5pt]
         $\bullet$\; \texttt{ALL previous\_results} {\scriptsize (str)} ---
          full accumulated output from Steps 1--3
          (omics + KG + literature)\\[2pt]
         $\bullet$\; \texttt{top\_genes} {\scriptsize (list[str])} ---
          ranked DEGs \emph{from OmicMiningAgent}\\[2pt]
         $\bullet$\; \texttt{paper\_dois} {\scriptsize (list[str])} ---
          literature DOIs \emph{from PubMedResearcher},
          used for citation cross-referencing};

    \draw[flowline] (step.south) -- (reads.north);

    % === PROMPT HIGHLIGHTS ===
    \node[infopanel, below=0.4cm of reads] (prompt)
        {\textbf{System Prompt Highlights}
         {\scriptsize (from \texttt{ScientistsAgent.system\_message})}\\[5pt]
         \textbf{1.}\; ``Analyse the incoming research task and
          determine which expert(s) are best suited''\\[2pt]
         \textbf{2.}\; ``Delegate with specific, tailored instructions
          for each expert''\\[2pt]
         \textbf{3.}\; ``Synthesise their responses into a unified,
          high-quality scientific analysis''\\[2pt]
         \textbf{4.}\; ``Prefer experts with \textbf{(RAG KB loaded)}
          --- they can retrieve evidence from real papers''\\[2pt]
         \textbf{5.}\; ``ALWAYS include a bioinformatics perspective
          when data-analysis methodology is involved''\\[2pt]
         \textbf{6.}\; ``ALWAYS include a biostatistics perspective
          when statistical claims or ageing are involved''\\[2pt]
         \textbf{7.}\; ``For hypothesis-generation tasks, use ALL
          relevant experts for multi-perspective coverage''};

    \draw[flowline] (reads.south) -- (prompt.north);

    % === ROUTING + EXPERT PANEL ===
    \node[toolpanel, below=0.4cm of prompt] (routing)
        {\textbf{PI Router:}\;
         Structured LLM output via \texttt{ExpertRouting}
         Pydantic model $\rightarrow$ selects 2--3 experts with
         tailored sub-queries};

    \draw[flowline] (prompt.south) -- (routing.north);

    % --- 4 Expert cards in a row ---
    \node[expertcard, fill=yellow!8,
          below=0.5cm of routing, xshift=-7cm] (exp1)
        {\textbf{Genomics-}\\
         \textbf{Expert}\\[3pt]
         \colorbox{green!15}{\scriptsize RAG KB}\\[3pt]
         {\tiny gene regulation,}\\
         {\tiny variant interp.,}\\
         {\tiny CRISPRi/a/ko,}\\
         {\tiny pharmacogenomics}};

    \node[expertcard, fill=pink!8,
          right=0.5cm of exp1] (exp2)
        {\textbf{Neuroscience-}\\
         \textbf{Expert}\\[3pt]
         \colorbox{green!15}{\scriptsize RAG KB}\\[3pt]
         {\tiny A$\beta$, tau, TREM2,}\\
         {\tiny microglia DAM,}\\
         {\tiny neuroinflammation,}\\
         {\tiny CSF biomarkers}};

    \node[expertcard, fill=lime!8,
          right=0.5cm of exp2] (exp3)
        {\textbf{Longevity/}\\
         \textbf{BiostatsExpert}\\[3pt]
         \colorbox{green!15}{\scriptsize RAG KB}\\[3pt]
         {\tiny ageing hallmarks,}\\
         {\tiny epigenetic clocks,}\\
         {\tiny mTOR/AMPK/sirtuins,}\\
         {\tiny survival analysis}};

    \node[expertcard, fill=cyan!8,
          right=0.5cm of exp3] (exp4)
        {\textbf{Cancer-}\\
         \textbf{Expert}\\[3pt]
         \colorbox{red!15}{\scriptsize LLM-only}\\[3pt]
         {\tiny tumor microenv.,}\\
         {\tiny oncogenesis,}\\
         {\tiny cancer therapeutics,}\\
         {\tiny resistance mech.}};

    % Fan-out arrows from routing to experts
    \draw[fanline] (routing.south) -- ++(0,-0.25) -| (exp1.north);
    \draw[fanline] (routing.south) -- ++(0,-0.25) -| (exp2.north);
    \draw[fanline] (routing.south) -- ++(0,-0.25) -| (exp3.north);
    \draw[fanline] (routing.south) -- ++(0,-0.25) -| (exp4.north);

    % Shared tool below experts
    \node[toolpanel, minimum width=19cm, text width=18cm,
          below=0.5cm of $(exp2.south)!0.5!(exp3.south)$] (ragtool)
        {\textbf{Shared Tool:}\;
         \texttt{scientist\_rag\_tool(query, expert)}\\[3pt]
         {\scriptsize Queries the specified expert's HyperRAG knowledge
          base built from curated author publications.\;
          Uses hypergraph-based retrieval with OpenAI embeddings and
          GPT-4.1 for answer generation.\;
          \textbf{RAG~KB} experts retrieve evidence from real papers;\;
          \textbf{LLM-only} experts reason without a private KB when
          their author KB has not yet been built.}\\[3pt]
         {\scriptsize \textit{See companion illustration for Scientist
          RAG initialisation and retrieval pipeline.}}};

    % Fan-in to shared tool
    \draw[fanline] (exp1.south) -- ++(0,-0.2) -| (ragtool.north west);
    \draw[fanline] (exp2.south) -- ++(0,-0.2) -|
                   ($(ragtool.north)!0.33!(ragtool.north west)$);
    \draw[fanline] (exp3.south) -- ++(0,-0.2) -|
                   ($(ragtool.north)!0.33!(ragtool.north east)$);
    \draw[fanline] (exp4.south) -- ++(0,-0.2) -| (ragtool.north east);

    % === PI SYNTHESIS ===
    \node[infopanel, draw=teal!50, fill=teal!5,
          below=0.4cm of ragtool, minimum width=20cm,
          text width=18.5cm] (synth)
        {\iconRobot{teal!60!black}\;\;
         \textbf{PI Synthesis}\\[3pt]
         {\scriptsize Integrates expert perspectives into unified
          analysis.\; Ranks mechanistic hypotheses by evidence
          strength.\; Proposes specific validation experiments
          per hypothesis.}};

    \draw[flowline] (ragtool.south) -- (synth.north);

    % === WRITES ===
    \node[rwpanel, draw=blue!30!black, fill=stateBlue,
          below=0.4cm of synth, minimum width=20cm,
          text width=18.5cm] (writes)
        {\textbf{$\blacktriangle$ Writes to AgentState}\\[5pt]
         $\bullet$\; \texttt{agent\_outputs[]} {\scriptsize (str)} ---
          multi-expert synthesis including:\\[2pt]
         \hphantom{$\bullet$\;}
          {\scriptsize --- ranked mechanistic hypotheses with
           Gene$\rightarrow$Pathway$\rightarrow$Phenotype chains}\\[1pt]
         \hphantom{$\bullet$\;}
          {\scriptsize --- per-expert domain perspectives and
           evidence summaries}\\[1pt]
         \hphantom{$\bullet$\;}
          {\scriptsize --- proposed validation experiments with
           quantitative readouts}\\[3pt]
         $\bullet$\; \texttt{expert\_routing\{\}} {\scriptsize (dict)} ---
          which experts were selected, why, and their tailored
          sub-queries};

    \draw[flowline] (synth.south) -- (writes.north);

    % === DOWNSTREAM CONSUMERS ===
    \node[consumerlabel, below=0.3cm of writes, text width=20cm]
          (consumers)
        {$\downarrow$\; Consumed by:\;
         \textbf{GoogleSearcher}
          (hypotheses to validate translationally)\; $\cdot$\;
         \textbf{ReportingNode}
          (hypothesis section, validation plans, expert routing log)};

\end{tikzpicture}%
}% end resizebox
\caption{\textbf{ScientistsAgent} --- workflow interface view.
The heaviest reader of \texttt{AgentState}, this compound agent
consumes all prior omics, KG, and literature output.  A PI-level
router uses a structured \texttt{ExpertRouting} Pydantic model to
select 2--3 of the four domain experts configured in
\texttt{configs/experts.yaml} and dispatch tailored sub-queries.
Genomics, Neuroscience, and Longevity/Biostatistics experts are
backed by curated-author HyperRAG KBs (\emph{RAG~KB}); CancerExpert
is configured but its KB is not yet built and is therefore routed
\emph{LLM-only}.  The PI synthesis node integrates expert
perspectives into ranked hypotheses with proposed validation
experiments.}
\label{fig:closeup-scientists}
\end{figure}

% ============================================================
\begin{figure}[htbp]
\centering
\resizebox{.78\textwidth}{!}{%
\begin{tikzpicture}[node distance=0.35cm, font=\sffamily,
    stepbar/.style={rectangle, rounded corners=6pt, thick,
        minimum width=20cm, align=center, font=\sffamily,
        drop shadow={opacity=0.10}, inner sep=6pt, text width=19cm},
    rwpanel/.style={rectangle, rounded corners=5pt, thick,
        minimum width=20cm, align=left, inner sep=6pt,
        font=\sffamily\small, text width=18.5cm},
    infopanel/.style={rectangle, rounded corners=5pt, draw=gray!50, thick,
        minimum width=20cm, align=left, fill=white,
        inner sep=6pt, font=\sffamily\small, text width=18.5cm},
    toolpanel/.style={rectangle, rounded corners=5pt, draw=gray!60, thick,
        minimum width=20cm, align=left, fill=gray!4,
        inner sep=6pt, font=\sffamily\small, text width=18.5cm},
    expertcard/.style={rectangle, rounded corners=4pt, draw=gray!50,
        thick, minimum width=4.2cm, minimum height=2.2cm,
        align=center, font=\sffamily\scriptsize, inner sep=5pt},
    flowline/.style={->, >=Stealth, very thick, draw=gray!40},
    fanline/.style={->, >=Stealth, thick, draw=teal!40},
    consumerlabel/.style={font=\sffamily\small\itshape, text=gray!60!black,
        align=center, text width=20cm}
]
    % === STEP BAR ===
    \node[stepbar, draw=teal!50!black, fill=agentTeal] (step)
        {\iconBulb{teal!60!black}\;\;
         \textbf{\large ScientistsAgent --- Multi-Expert Hypothesis
          Synthesis}\\[3pt]
         {\small Planner: \textit{``Synthesise mechanisms via expert
          panel, score hypotheses, propose experiments.''}}\\[2pt]
         {\scriptsize\sffamily\color{gray!60!black}
          Shown as a fourth step here; the planner dynamically
          decides invocation order and may loop back to earlier
          agents after synthesis reveals knowledge gaps.}};

    % === READS ===
    \node[rwpanel, draw=green!40!black, fill=agentGreen,
          below=0.4cm of step] (reads)
        {\textbf{$\blacktriangledown$ Reads from AgentState}
         {\scriptsize (heaviest reader --- consumes all prior output)}\\[5pt]
         $\bullet$\; \texttt{ALL previous\_results} {\scriptsize (str)} ---
          full accumulated output from Steps 1--3
          (omics + KG + literature)\\[2pt]
         $\bullet$\; \texttt{top\_genes} {\scriptsize (list[str])} ---
          ranked DEGs \emph{from OmicMiningAgent}\\[2pt]
         $\bullet$\; \texttt{paper\_dois} {\scriptsize (list[str])} ---
          literature DOIs \emph{from PubMedResearcher},
          used for citation cross-referencing};

    \draw[flowline] (step.south) -- (reads.north);

    % === PROMPT HIGHLIGHTS ===
    \node[infopanel, below=0.4cm of reads] (prompt)
        {\textbf{System Prompt Highlights}
         {\scriptsize (from \texttt{ScientistsAgent.system\_message})}\\[5pt]
         \textbf{1.}\; ``Analyse the incoming research task and
          determine which expert(s) are best suited''\\[2pt]
         \textbf{2.}\; ``Delegate with specific, tailored instructions
          for each expert''\\[2pt]
         \textbf{3.}\; ``Synthesise their responses into a unified,
          high-quality scientific analysis''\\[2pt]
         \textbf{4.}\; ``Prefer experts with \textbf{(RAG KB loaded)}
          --- they can retrieve evidence from real papers''\\[2pt]
         \textbf{5.}\; ``ALWAYS include a bioinformatics perspective
          when data-analysis methodology is involved''\\[2pt]
         \textbf{6.}\; ``ALWAYS include a biostatistics perspective
          when statistical claims or ageing are involved''\\[2pt]
         \textbf{7.}\; ``For hypothesis-generation tasks, use ALL
          relevant experts for multi-perspective coverage''};

    \draw[flowline] (reads.south) -- (prompt.north);

    % === ROUTING + EXPERT PANEL ===
    \node[toolpanel, below=0.4cm of prompt] (routing)
        {\textbf{PI Router:}\;
         Structured LLM output via \texttt{ExpertRouting}
         Pydantic model $\rightarrow$ selects 2--3 experts with
         tailored sub-queries};

    \draw[flowline] (prompt.south) -- (routing.north);

    % --- 4 Expert cards in a row ---
    \node[expertcard, fill=yellow!8,
          below=0.5cm of routing, xshift=-7cm] (exp1)
        {\textbf{Genomics-}\\
         \textbf{Expert}\\[3pt]
         \colorbox{green!15}{\scriptsize RAG KB}\\[3pt]
         {\tiny gene regulation,}\\
         {\tiny variant interp.,}\\
         {\tiny CRISPRi/a/ko,}\\
         {\tiny pharmacogenomics}};

    \node[expertcard, fill=pink!8,
          right=0.5cm of exp1] (exp2)
        {\textbf{Neuroscience-}\\
         \textbf{Expert}\\[3pt]
         \colorbox{green!15}{\scriptsize RAG KB}\\[3pt]
         {\tiny A$\beta$, tau, TREM2,}\\
         {\tiny microglia DAM,}\\
         {\tiny neuroinflammation,}\\
         {\tiny CSF biomarkers}};

    \node[expertcard, fill=lime!8,
          right=0.5cm of exp2] (exp3)
        {\textbf{Longevity/}\\
         \textbf{BiostatsExpert}\\[3pt]
         \colorbox{green!15}{\scriptsize RAG KB}\\[3pt]
         {\tiny ageing hallmarks,}\\
         {\tiny epigenetic clocks,}\\
         {\tiny mTOR/AMPK/sirtuins,}\\
         {\tiny survival analysis}};

    \node[expertcard, fill=cyan!8,
          right=0.5cm of exp3] (exp4)
        {\textbf{Cancer-}\\
         \textbf{Expert}\\[3pt]
         \colorbox{red!15}{\scriptsize LLM-only}\\[3pt]
         {\tiny tumor microenv.,}\\
         {\tiny oncogenesis,}\\
         {\tiny cancer therapeutics,}\\
         {\tiny resistance mech.}};

    % Fan-out arrows from routing to experts
    \draw[fanline] (routing.south) -- ++(0,-0.25) -| (exp1.north);
    \draw[fanline] (routing.south) -- ++(0,-0.25) -| (exp2.north);
    \draw[fanline] (routing.south) -- ++(0,-0.25) -| (exp3.north);
    \draw[fanline] (routing.south) -- ++(0,-0.25) -| (exp4.north);

    % Shared tool below experts
    \node[toolpanel, minimum width=19cm, text width=18cm,
          below=0.5cm of $(exp2.south)!0.5!(exp3.south)$] (ragtool)
        {\textbf{Shared Tool:}\;
         \texttt{scientist\_rag\_tool(query, expert)}\\[3pt]
         {\scriptsize Queries the specified expert's HyperRAG knowledge
          base built from curated author publications.\;
          Uses hypergraph-based retrieval with OpenAI embeddings and
          GPT-4.1 for answer generation.\;
          \textbf{RAG~KB} experts retrieve evidence from real papers;\;
          \textbf{LLM-only} experts reason without a private KB when
          their author KB has not yet been built.}\\[3pt]
         {\scriptsize \textit{See companion illustration for Scientist
          RAG initialisation and retrieval pipeline.}}};

    % Fan-in to shared tool
    \draw[fanline] (exp1.south) -- ++(0,-0.2) -| (ragtool.north west);
    \draw[fanline] (exp2.south) -- ++(0,-0.2) -|
                   ($(ragtool.north)!0.33!(ragtool.north west)$);
    \draw[fanline] (exp3.south) -- ++(0,-0.2) -|
                   ($(ragtool.north)!0.33!(ragtool.north east)$);
    \draw[fanline] (exp4.south) -- ++(0,-0.2) -| (ragtool.north east);

    % === PI SYNTHESIS ===
    \node[infopanel, draw=teal!50, fill=teal!5,
          below=0.4cm of ragtool, minimum width=20cm,
          text width=18.5cm] (synth)
        {\iconRobot{teal!60!black}\;\;
         \textbf{PI Synthesis}\\[3pt]
         {\scriptsize Integrates expert perspectives into unified
          analysis.\; Ranks mechanistic hypotheses by evidence
          strength.\; Proposes specific validation experiments
          per hypothesis.}};

    \draw[flowline] (ragtool.south) -- (synth.north);

    % === WRITES ===
    \node[rwpanel, draw=blue!30!black, fill=stateBlue,
          below=0.4cm of synth, minimum width=20cm,
          text width=18.5cm] (writes)
        {\textbf{$\blacktriangle$ Writes to AgentState}\\[5pt]
         $\bullet$\; \texttt{agent\_outputs[]} {\scriptsize (str)} ---
          multi-expert synthesis including:\\[2pt]
         \hphantom{$\bullet$\;}
          {\scriptsize --- ranked mechanistic hypotheses with
           Gene$\rightarrow$Pathway$\rightarrow$Phenotype chains}\\[1pt]
         \hphantom{$\bullet$\;}
          {\scriptsize --- per-expert domain perspectives and
           evidence summaries}\\[1pt]
         \hphantom{$\bullet$\;}
          {\scriptsize --- proposed validation experiments with
           quantitative readouts}\\[3pt]
         $\bullet$\; \texttt{expert\_routing\{\}} {\scriptsize (dict)} ---
          which experts were selected, why, and their tailored
          sub-queries};

    \draw[flowline] (synth.south) -- (writes.north);

    % === DOWNSTREAM CONSUMERS ===
    \node[consumerlabel, below=0.3cm of writes, text width=20cm]
          (consumers)
        {$\downarrow$\; Consumed by:\;
         \textbf{GoogleSearcher}
          (hypotheses to validate translationally)\; $\cdot$\;
         \textbf{ReportingNode}
          (hypothesis section, validation plans, expert routing log)};

\end{tikzpicture}%
}% end resizebox
\caption{\textbf{ScientistsAgent} --- workflow interface view.
The heaviest reader of \texttt{AgentState}, this compound agent
consumes all prior omics, KG, and literature output.  A PI-level
router uses a structured \texttt{ExpertRouting} Pydantic model to
select 2--3 of the four domain experts configured in
\texttt{configs/experts.yaml} and dispatch tailored sub-queries.
Genomics, Neuroscience, and Longevity/Biostatistics experts are
backed by curated-author HyperRAG KBs (\emph{RAG~KB}); CancerExpert
is configured but its KB is not yet built and is therefore routed
\emph{LLM-only}.  The PI synthesis node integrates expert
perspectives into ranked hypotheses with proposed validation
experiments.}
\label{fig:closeup-scientists}
\end{figure}
